\appendix
\section{Proof of Mean-Dominance for Extreme Entries}
\label{app:outlier_bias_proof}

This appendix proves a formal version of the empirical observation in
Section~\ref{subsec:outlier_mean_dominance}:
among the largest-magnitude entries of the activation matrix, the broadcasted
mean component dominates.

\subsection{Setup and notation}

Let $x_i \in \mathbb{R}^H$ denote token activations and let $X\in\mathbb{R}^{N\times H}$
stack them row-wise. We model each entry as
\begin{equation}
\label{eq:app_model_entry}
x_{ij} = \mu_j + s_{ij},
\end{equation}
where $\mu\in\mathbb{R}^H$ is deterministic (the coherent mean) and
$\{s_{ij}\}$ are random fluctuations with $\mathbb{E}[s_{ij}]=0$.

Define the mean matrix $M\triangleq \mathbf{1}\mu^\top$ and the centered matrix
$S\triangleq X-M$ so that $S_{ij}=s_{ij}$.

For any $\varepsilon\in(0,1)$, define the (population) outlier set as the indices
of the largest $\varepsilon NH$ entries by magnitude:
\begin{equation}
\label{eq:app_outlier_set}
\Omega_\varepsilon
\;\triangleq\;
\left\{(i,j): |x_{ij}| \text{ is among the top } \varepsilon NH \text{ values of } \{|x_{i'j'}|\}_{i',j'} \right\}.
\end{equation}
We will show that under mild conditions, entries in $\Omega_\varepsilon$ have
$s_{ij}$ small relative to $\mu_j$, hence the mean dominates.

\subsection{Assumptions}

\begin{assumption}[Sub-Gaussian fluctuations]
\label{assump:subg}
The entries $\{s_{ij}\}$ are independent and sub-Gaussian with parameter $\sigma^2$:
for all $t\ge 0$,
\[
\mathbb{P}\big(|s_{ij}| \ge t\big) \le 2\exp\!\left(-\frac{t^2}{2\sigma^2}\right).
\]
\end{assumption}

\begin{assumption}[Mean-gap and mass near the maximum]
\label{assump:gap}
Let $m \triangleq \max_{1\le j\le H} |\mu_j|$.
Fix $\gamma\in(0,1)$ and define the ``high-mean'' coordinate set
\[
\mathcal{J}_\gamma \triangleq \{j: |\mu_j|\ge (1-\gamma)m\},
\qquad
p_\gamma \triangleq |\mathcal{J}_\gamma|/H.
\]
Assume $p_\gamma>0$ and that for $j\notin\mathcal{J}_\gamma$, we have the strict gap
$|\mu_j|\le (1-\gamma)m$.
\end{assumption}

\begin{assumption}[Spike-dominance regime]
\label{assump:regime}
The mean amplitude satisfies
\[
m \;\ge\; \frac{4\sigma}{\gamma}\sqrt{2\log\!\left(\frac{4}{\varepsilon p_\gamma}\right)}.
\]
\end{assumption}

Assumption~\ref{assump:gap} captures the empirically observed phenomenon that a
nontrivial fraction of coordinates have mean values close to the maximum mean.
Assumption~\ref{assump:regime} places us in a regime where the mean dominates
rare noise events at the outlier quantile.

\subsection{Main theorem}

\begin{theorem}[Mean dominates the top-$\varepsilon$ fraction of entries]
\label{thm:mean_dominates_outliers}
Under Assumptions~\ref{assump:subg}--\ref{assump:regime}, with probability at least
\[
1 - 2\exp\!\left(-\frac{\varepsilon p_\gamma NH}{32}\right) - 2\varepsilon NH \exp\!\left(-\frac{\gamma^2 m^2}{32\sigma^2}\right),
\]
the following two statements hold:

\begin{enumerate}
\item (\emph{Outliers occur on high-mean coordinates.})
Every outlier index $(i,j)\in \Omega_\varepsilon$ satisfies $j\in \mathcal{J}_\gamma$
and
\[
|s_{ij}| \le \frac{\gamma m}{4}.
\]

\item (\emph{Per-outlier mean dominance.})
For every $(i,j)\in \Omega_\varepsilon$,
the squared-magnitude share of the mean satisfies
\[
\frac{M_{ij}^2}{M_{ij}^2 + S_{ij}^2}
\;\ge\;
1 - \left(\frac{\gamma}{4(1-\gamma)}\right)^2.
\]
\end{enumerate}
\end{theorem}

\subsection{Proof}

We prove two counting lemmas: (i) there are enough ``clean'' entries on high-mean
coordinates to fill the outlier set, and (ii) low-mean coordinates cannot produce
enough entries above the outlier threshold without requiring rare large noise.

\begin{lemma}[Many high-mean entries have small noise]
\label{lem:many_clean_highmean}
Let $t \triangleq \gamma m/4$.
Then with probability at least $1-2\exp(-\varepsilon p_\gamma NH/32)$, the number of
indices $(i,j)$ with $j\in\mathcal{J}_\gamma$ and $|s_{ij}|\le t$ is at least
\[
\frac{1}{2}p_\gamma NH.
\]
\end{lemma}

\begin{proof}
For each $(i,j)$ with $j\in\mathcal{J}_\gamma$, define the indicator
$Z_{ij} \triangleq \mathbf{1}\{|s_{ij}|\le t\}$.
By Assumption~\ref{assump:subg},
\[
\mathbb{P}(Z_{ij}=1) = 1-\mathbb{P}(|s_{ij}|>t)
\;\ge\;
1-2\exp\!\left(-\frac{t^2}{2\sigma^2}\right).
\]
Under Assumption~\ref{assump:regime}, $t=\gamma m/4$ implies
$2\exp(-t^2/(2\sigma^2)) \le \varepsilon p_\gamma/2$,
hence $\mathbb{P}(Z_{ij}=1)\ge 1-\varepsilon p_\gamma/2\ge 1/2$.
The total count is
\[
Z \triangleq \sum_{j\in\mathcal{J}_\gamma}\sum_{i=1}^N Z_{ij},
\qquad
\mathbb{E}[Z] \ge \frac{1}{2} p_\gamma NH.
\]
Since the $Z_{ij}$ are independent Bernoulli variables, a Chernoff bound yields
\[
\mathbb{P}\!\left(Z \le \frac{1}{2}\mathbb{E}[Z]\right)
\le
\exp\!\left(-\frac{\mathbb{E}[Z]}{8}\right)
\le
\exp\!\left(-\frac{p_\gamma NH}{32}\right)
\le
\exp\!\left(-\frac{\varepsilon p_\gamma NH}{32}\right),
\]
which proves the claim (up to a constant-factor tightening; we use the displayed bound).
\end{proof}

\begin{lemma}[Low-mean coordinates rarely exceed the high-mean threshold]
\label{lem:lowmean_rare}
Let $t=\gamma m/4$ and define the deterministic threshold
\[
\tau \triangleq (1-\gamma)m - t.
\]
For any $j\notin\mathcal{J}_\gamma$ and any $i$,
\[
\mathbb{P}\big(|x_{ij}|\ge \tau\big)
\le
2\exp\!\left(-\frac{t^2}{2\sigma^2}\right).
\]
Consequently, with probability at least $1-2\varepsilon NH\exp(-t^2/(2\sigma^2))$,
the number of indices $(i,j)$ with $j\notin\mathcal{J}_\gamma$ and $|x_{ij}|\ge\tau$
is at most $\varepsilon NH$.
\end{lemma}

\begin{proof}
Fix $j\notin\mathcal{J}_\gamma$, so $|\mu_j|\le (1-\gamma)m$ by Assumption~\ref{assump:gap}.
If $|x_{ij}|=|\mu_j+s_{ij}|\ge \tau=(1-\gamma)m-t$, then necessarily
$|s_{ij}|\ge t$ (otherwise $|\mu_j+s_{ij}|\le |\mu_j|+|s_{ij}|\le (1-\gamma)m+t < \tau$).
Thus
\[
\mathbb{P}(|x_{ij}|\ge\tau)\le \mathbb{P}(|s_{ij}|\ge t)
\le 2\exp\!\left(-\frac{t^2}{2\sigma^2}\right).
\]
The second statement follows by a union bound over all such indices and Markov/Chernoff-style counting; we use a simple union bound:
\[
\mathbb{P}\left(\exists\,\varepsilon NH \text{ distinct indices with } j\notin\mathcal{J}_\gamma,\,|x_{ij}|\ge\tau\right)
\le
2\varepsilon NH\exp\!\left(-\frac{t^2}{2\sigma^2}\right).
\]
\end{proof}

\begin{proof}[Proof of Theorem~\ref{thm:mean_dominates_outliers}]
Let $t=\gamma m/4$ and $\tau=(1-\gamma)m-t$ as above.

By Lemma~\ref{lem:many_clean_highmean}, with high probability there exist at least
$\frac{1}{2}p_\gamma NH$ indices $(i,j)$ with $j\in\mathcal{J}_\gamma$ and
$|s_{ij}|\le t$.
For each such index, we have
\[
|x_{ij}| = |\mu_j+s_{ij}|
\ge |\mu_j|-|s_{ij}|
\ge (1-\gamma)m - t
= \tau.
\]
Therefore, there are at least $\frac{1}{2}p_\gamma NH$ entries with magnitude at least $\tau$.

Under Assumption~\ref{assump:regime}, $p_\gamma/2 \ge \varepsilon$ (equivalently, $p_\gamma \ge 2\varepsilon$),
so these entries alone exceed the required $\varepsilon NH$ outliers.
Hence the outlier set $\Omega_\varepsilon$ is contained within the set of indices
with $|x_{ij}|\ge\tau$.

Now apply Lemma~\ref{lem:lowmean_rare}:
with high probability, at most $\varepsilon NH$ entries from low-mean coordinates
($j\notin\mathcal{J}_\gamma$) can reach magnitude $\tau$.
Since $\Omega_\varepsilon$ consists of the top $\varepsilon NH$ entries,
it follows that, with the stated probability, every $(i,j)\in\Omega_\varepsilon$
must satisfy $j\in\mathcal{J}_\gamma$ and $|s_{ij}|\le t$.
This proves Item (1).

For Item (2), for $(i,j)\in\Omega_\varepsilon$ we have $|\mu_j|\ge (1-\gamma)m$ and
$|s_{ij}|\le t=\gamma m/4$, hence
\[
\frac{M_{ij}^2}{M_{ij}^2 + S_{ij}^2}
=
\frac{\mu_j^2}{\mu_j^2 + s_{ij}^2}
\ge
\frac{(1-\gamma)^2m^2}{(1-\gamma)^2m^2 + (\gamma m/4)^2}
=
1 - \left(\frac{\gamma}{4(1-\gamma)}\right)^2,
\]
as required.
\end{proof}

\paragraph{Remark.}
The theorem formalizes a qualitative mechanism: when a nontrivial fraction of coordinates
carry large coherent means and the noise is sub-Gaussian, the largest-magnitude entries
come primarily from those mean-heavy coordinates, so mean subtraction removes the dominant
source of outliers. The residual spike--tail structure then reflects variance anisotropy,
consistent with Section~\ref{subsec:spectral_role_bias}.